%% Commands for TeXCount
%TC:macro \cite [option:text,text]
%TC:macro \citep [option:text,text]
%TC:macro \citet [option:text,text]
%TC:envir table 0 1
%TC:envir table* 0 1
%TC:envir tabular [ignore] word
%TC:envir displaymath 0 word
%TC:envir math 0 word
%TC:envir comment 0 0
%%
%%
%% For submission and review:
%% \documentclass[acmsmall,screen,review,anonymous]{acmart}.

\documentclass[acmsmall,screen,nonacm]{acmart}

% \setcopyright{rightsretained}
% \acmDOI{10.1145/3674642}
% \acmYear{2024}
% \acmJournal{PACMPL}
% \acmVolume{8}
% \acmNumber{ICFP}
% \acmArticle{253}
% \acmMonth{8}
% \acmSubmissionID{icfp24main-p58-p}
% \received{2024-02-28}
% \received[accepted]{2024-06-18}

\citestyle{acmnumeric}

\usepackage{mathpartir}
\usepackage{wrapfig}
\usepackage{scalerel}
\usepackage{minted}
\usepackage{mathtools}
\usepackage{fontawesome} % lock symbol
\usepackage{tikz}
\usepackage{stmaryrd}
\usepackage[bb=boondox]{mathalfa}
\usepackage{pifont}

\newcommand{\under}[2]{\underset{\raisebox{2mm}{$\scriptstyle #2$}}{#1}}
\newcommand{\inh}[1]{\under{#1}{\uparrow}}
\newcommand{\syn}[1]{\under{#1}{\downarrow}}
\DeclareMathOperator{\llet}{\text{let}}
\DeclareMathOperator{\lin}{\text{in}}
\DeclareMathOperator{\case}{\text{case}}
\DeclareMathOperator{\bbox}{\text{box}}
\DeclareMathOperator{\unbox}{\text{unbox}}
\DeclareMathOperator{\mmod}{\text{mod}}
\DeclareMathOperator{\thunk}{\text{thunk}}
\DeclareMathOperator{\force}{\text{force}}
\DeclareMathOperator{\return}{\text{return}}
\newcommand{\unmod}[3]{\text{let}_{#1} \text{ mod}_{#2}({#3}) \leftarrow}
\newcommand{\at}{\mathbin{\mbox{\small\ensuremath{@}}}}
\DeclarePairedDelimiter\ann{\langle}{\rangle}
\newcommand{\unit}{\mathbb{1}}
\newcommand{\zero}{\mathbb{0}}
\DeclareMathOperator{\inl}{\text{inl}}
\DeclareMathOperator{\inr}{\text{inr}}
\DeclareMathOperator{\suc}{\text{succ}}
\DeclareMathOperator{\absurd}{\text{absurd}}
\DeclareMathOperator{\primrec}{\text{primrec}}
\newcommand{\nat}{\mathbb{N}}
\newcommand{\tee}{\vdash}
\newcommand{\teev}{\vdash^\mathrm{v}}
\newcommand{\teec}{\vdash^\mathrm{c}}
\newcommand{\hole}{\,?}
\newcommand*{\Ev}{E_\mathrm{v}}
\newcommand*{\Ep}{K}
\DeclareMathOperator{\red}{\rightsquigarrow}
\DeclarePairedDelimiter\floor{\lfloor}{\rfloor}
\DeclarePairedDelimiter\ceil{\lceil}{\rceil}
\DeclareMathOperator{\lock}{\tikz{\node[scale=1, line width=0pt, inner sep=0pt] at (0,0) {\faLock\,};}}
\definecolor{gainsboro}{rgb}{0.86, 0.86, 0.86}
\newcommand{\reducedstrut}{\vrule width 0pt height .9\ht\strutbox depth .9\dp\strutbox\relax}
\newcommand{\gray}[1]{%
  \begingroup
  \setlength{\fboxsep}{0pt}%
  \colorbox{gainsboro}{\reducedstrut$#1$\/}%
  \endgroup
}

\newcommand{\clashes}{\,\text{\Lightning}\,}
\newcommand{\notclashes}{\,{\text{\Lightning}\mkern-8.5mu\backslash}\,}
\DeclareMathOperator{\joinctx}{\mathbin{\fatsemi}}
\DeclareMathOperator{\splitctx}{\curlyveeuparrow}

\newcommand{\one}{\mathbf{1}}
\newcommand{\sem}[1]{\llbracket #1 \rrbracket}

\renewcommand{\qed}{\hfill\resizebox{8px}{8px}{\input{image.pdf_tex}}}

\newtheorem{law}{Law}

\begin{document}

\title{Internship Report: Call-By-Push-Value Calculus}

\author{Isabelle Gupta}
\affiliation{%
  \institution{University of Edinburgh}
  %\city{Edinburgh}
  \country{United Kingdom}
}

% \author{Anton Lorenzen}
% \orcid{0000-0003-3538-9688}
% \affiliation{%
  % \institution{University of Edinburgh}
  % %\city{Edinburgh}
  % \country{United Kingdom}
% }
% \email{anton.lorenzen@ed.ac.uk}

% \author{Sam Lindley}
% \orcid{0000-0002-1360-4714}
% \affiliation{%
  % \institution{University of Edinburgh}
  % %\city{Edinburgh}
  % \country{United Kingdom}
% }
% \email{Sam.Lindley@ed.ac.uk}

%% By default, the full list of authors will be used in the page
%% headers. Often, this list is too long, and will overlap
%% other information printed in the page headers. This command allows
%% the author to define a more concise list
%% of authors' names for this purpose.
% \renewcommand{\shortauthors}{}

% \begin{abstract}
% \end{abstract}

%% The code below is generated by the tool at http://dl.acm.org/ccs.cfm.
% \begin{CCSXML}
% <ccs2012>
   % <concept>
       % <concept_id>10003752.10003790.10003793</concept_id>
       % <concept_desc>Theory of computation~Modal and temporal logics</concept_desc>
       % <concept_significance>500</concept_significance>
       % </concept>
   % <concept>
       % <concept_id>10003752.10003790.10003801</concept_id>
       % <concept_desc>Theory of computation~Linear logic</concept_desc>
       % <concept_significance>300</concept_significance>
       % </concept>
 % </ccs2012>
% \end{CCSXML}

% \ccsdesc[500]{Theory of computation~Modal and temporal logics}
% \ccsdesc[300]{Theory of computation~Linear logic}

%% Keywords. The author(s) should pick words that accurately describe
%% the work being presented. Separate the keywords with commas.
% \keywords{Modal Types, Mode Qualifiers}

\maketitle

\section{Graded Call-By-Push-Value}

A new calculus to fix the issues we identified with the graded modal type theory.
The idea is that we have a new value judgement $\Gamma \tee v :^q A$,
where, informally, $q\cdot\Gamma \tee v :^q A$ if and only if
$\Gamma \tee v : A$ in the graded modal type theory.

\paragraph{Syntax}

\begin{align*}
  A &::= U\, \underline{B} \mid \boxempty^q A \mid A \times A' \mid A + A' \mid \unit \\
  \underline{B} &::= F^q\, A \mid \,^q A \rightarrow \underline{B} \\
  v &::= x \mid \thunk\, e \mid \bbox_q\, v \mid (v_1, v_2) \mid () \mid \inl\, v \mid \inr\, v \\
  e &::= \return\, v \mid \lambda x.\, e \mid e\, v
      \mid \llet x \leftarrow e_1 \lin e_2 \\
    &\mid \llet_{q_1} \bbox_{q_2} x = v \lin e \\
    &\mid \llet_q (x, y) = v \lin e
      \mid \case_q\, v\, \{x.e_1;\; y.e_2\} \\
    &\mid \absurd \, e
\end{align*}

\paragraph{Value typing rules} \hfill \framebox{$\Gamma \tee v : A$}

\begin{mathpar}
  \inferrule*[right=var]{
    \phantom{.}
  }{
    0 \cdot \Gamma, x :^q A, 0 \cdot \Gamma' \tee x :^q A
  } \and
  \inferrule*[right=vsub]{
    q' \leq q'' \\
    \Gamma, x :^{q'} A, \Gamma' \tee v :^q B
  }{
    \Gamma, x :^{q''} A, \Gamma' \tee v :^q B
  } \and
  \inferrule*[right=thunk]{
    \Gamma \tee e : \underline{B}
  }{
    q\cdot\Gamma \tee \thunk\, e :^{q} U\, \underline{B}
  } \and
  \inferrule*[right=unit]{
    \phantom{.}
  }{
    0\cdot\Gamma \tee () :^q \unit
  } \and
  \inferrule*[right=pair]{
    \Gamma_1 \tee v_1 :^q A \\
    \Gamma_2 \tee v_2 :^q B
  }{
    \Gamma_1 + \Gamma_2 \tee (v_1, v_2) :^q A \times B
  } \and
  \inferrule*[right=inl]{
    \Gamma \tee v :^q A
  }{
    \Gamma \tee \inl\, v :^q A + B
  } \and
  \inferrule*[right=inr]{
    \Gamma \tee v :^q B
  }{
    \Gamma \tee \inr\, v :^q A + B
  } \and
  \inferrule*[right=$\boxempty$I]{
    \Gamma \tee v :^{q'q} A
  }{
    \Gamma \tee \bbox_q\, v :^{q'} \boxempty^q A
  }
\end{mathpar}

\paragraph{Expression typing rules} \hfill \framebox{$\Gamma \tee e : \underline{B}$}

\begin{mathpar}
  \inferrule*[right=return]{
    \Gamma \tee v :^q A
  }{
    \Gamma \tee \return\, v : F^q\, A
  } \and
  \inferrule*[right=force]{
    \Gamma \tee v :^1 U\, \underline{B}
  }{
    \Gamma \tee \force\, v : \underline{B}
  } \and
  \inferrule*[right=esub]{
    q \leq q' \\
    \Gamma, x :^{q} A, \Gamma' \tee e : \underline{B}
  }{
    \Gamma, x :^{q'} A, \Gamma' \tee e : \underline{B}
  } \and
  \inferrule*[right=lam]{
    \Gamma, x :^q A \tee e : \underline{B}
  }{
    \Gamma \tee \lambda x.\, e : \,^q A \rightarrow \underline{B}
  } \and
  \inferrule*[right=app]{
    \Gamma_1 \tee e : \,^q A \rightarrow \underline{B} \\
    \Gamma_2 \tee v :^q A
  }{
    \Gamma_1 + \Gamma_2 \tee e\, v : \underline{B}
  } \and
  \inferrule*[right=bind]{
    \Gamma_1 \tee e_1 : F^q\, A \\
    \Gamma_2, x :^q A \tee e_2 : \underline{B}
  }{
    \Gamma_1 + \Gamma_2 \tee \llet x \leftarrow e_1 \lin e_2 : \underline{B}
  } \and
  \inferrule*[right=$\boxempty$E]{
    \Gamma_1 \tee v :^{q_1} \boxempty^{q_2} A \\
    \Gamma_2, x :^{q_1 \cdot q_2} A \tee e : \underline{B}
  }{
    \Gamma_1 + \Gamma_2 \tee \llet_{q_1} \bbox_{q_2} x = v \lin e : \underline{B}
  } \and
  \inferrule*[right=split]{
    \Gamma_1 \tee v :^{q} A_1 \times A_2 \\
    \Gamma_2, x :^q A_1, y :^q A_2 \tee e : \underline{B}
  }{
    \Gamma_1 + \Gamma_2 \tee \llet_q (x, y) = v \lin e : \underline{B}
  } \and
  \inferrule*[right=case]{
    \Gamma_1 \tee v :^q A_1 + A_2 \\
    \Gamma_2, x :^q A_1 \tee e_1 : \underline{B} \\
    \Gamma_2, y :^q A_2 \tee e_2 : \underline{B}
  }{
    \Gamma_1 + \Gamma_2 \tee \case_q\, v\, \{x.e_1;\; y.e_2\} : \underline{B}
  } \and
  \inferrule*[right=absurd]{
    \Gamma \tee e : F^q\,\zero
  }{
    \top\cdot\Gamma \tee \absurd \, e : \underline{B}
  }
\end{mathpar}

Parameterisation is as in the main report.

\section{Laws}

The laws that must hold for the algebra are:

\begin{law}
\label{law:add-assoc}
  $q_1 + (q_2 + q_3) = (q_1 + q_2) + q_3$.
\end{law}

\begin{law}
\label{law:add-comm}
  $q_1 + q_2 = q_2 + q_1$.
\end{law}

\begin{law}
\label{law:mul-assoc-1}
  $q_1 \cdot (q_2 \cdot q_3) \le (q_1 \cdot q_2) \cdot q_3$.
\end{law}

\begin{law}
\label{law:mul-distr-add-left-1}
  $q_1 \cdot q_2 + q_1 \cdot q_3 \le q_1 \cdot (q_2 + q_3)$.
\end{law}

\begin{law}
\label{law:mul-distr-add-right-1}
  $q_1 \cdot q_3 + q_2 \cdot q_3 \le (q_1 + q_2) \cdot q_3$.
\end{law}

\begin{law}
\label{law:mul-inverse-div}
  $q_2 \cdot (q_1 \setminus q_2) \le q_1$.
\end{law}

\begin{law}
\label{law:add-mono}
  If $q_1 \le q_3$ and $q_2 \le q_4$, then $q_1 + q_2 \le q_3 + q_4$.
\end{law}

\begin{law}
\label{law:mul-mono-left}
  If $q_2 \le q_3$, then $q_1 \cdot q_2 \le q_1 \cdot q_3$.
\end{law}

The laws that must hold for forward substitution are:

\begin{law}
\label{law:mul-assoc-2}
  $(q_1 \cdot q_2) \cdot q_3 \le q_1 \cdot (q_2 \cdot q_3)$.
  FOR BOX INTRO
\end{law}

\begin{law}
\label{law:mul-distr-add-left-2}
  $q_1 \cdot (q_2 + q_3) \le q_1 \cdot q_2 + q_1 \cdot q_3$.
\end{law}

\begin{law}
\label{law:mul-distr-add-right-2}
  $(q_1 + q_2) \cdot q_3 \le q_1 \cdot q_3 + q_2 \cdot q_3$.
  FOR BOX INTRO
\end{law}

\section{Proofs}

\begin{lemma}
\label{lem:mul-left-inverse-div-le}
  $q^\prime \le (q \cdot q^\prime) \setminus q$.
\end{lemma}

\begin{proof}
  Consider the set $M = \{q^{\prime\prime} : q \cdot q^{\prime\prime} \le q \cdot q^\prime\}$, the candidates for $q^{\prime\prime} = q \cdot q^\prime \setminus q$.

  Since $q \cdot q^\prime \le q \cdot q^\prime$, it follows that $q^\prime \in M$.
  
  Thus, since $q \cdot q^\prime \setminus q$ is given by $\bigvee M$, $q^\prime \le q \cdot q^\prime \setminus q$, as required.
  
\end{proof}

\begin{lemma}
\label{lem:div-mono}
  If $q_1 \le q_2$, then $q_1 \setminus q \le q_2 \setminus q$.
\end{lemma}

\begin{proof}
  Consider the sets $M_1 = \{q^\prime : q \cdot q^\prime \le q_1\}$ and $M_2 = \{q^\prime : q \cdot q^\prime = q_2\}$.

  Since $q_1 \le q_2$, every element in $M_1$ is in $M_2$, so $M_1 \subseteq M_2$.
  
  Thus, since $q_1 \setminus q$ and $q_2 \setminus q$ are given by $\bigvee M_1$ and $\bigvee M_2$, $q_1 \setminus q \le q_2 \setminus q$, as required.
  
\end{proof}

\begin{lemma}
\label{lem:div-distr-add}
  $q_1 \setminus q_3 + q_2 \setminus q_3 \le (q_1 + q_2) \setminus q_3$.
\end{lemma}

\begin{proof}
\end{proof}

\begin{lemma}
\label{lem:new-modal-mode-split}
  If $\Gamma \tee v :^q A$, then $q \cdot \Gamma^\prime \le \Gamma$ for some $\Gamma^\prime$, $q_1 \cdot \Gamma^\prime \tee v :^{q_1} A$ and $q_2 \cdot \Gamma^\prime \tee v :^{q_2} A$ for all $q_1 + q_2 = q^\prime$.
\end{lemma}

\begin{proof}
  Proceed by induction over the structure of the derivation of $\Gamma \tee v :^q A$.
  \begin{description}
    \item[Case VAR.]
    By assumption, have $0 \cdot \Gamma, x :^q A, 0 \cdot \Gamma^\prime \tee x :^q A$.
    
    For the exists statement, let $\Gamma^\prime = 0 \cdot \Gamma, x :^1 A, 0 \cdot \Gamma^\prime$.

    $0 \cdot \Gamma, x :^q A, 0 \cdot \Gamma^\prime \ge q \cdot (0 \cdot \Gamma, x :^1 A, 0 \cdot \Gamma^\prime)$, as required.

    $0 \cdot \Gamma, x :^{q_1} A, 0 \cdot \Gamma^\prime \tee x :^{q_1} A$ and $0 \cdot \Gamma, x :^{q_2} A, 0 \cdot \Gamma^\prime \tee x :^{q_2} A$, as required.
    \\
    \item[Case LAM.]
    By assumption, have \textbf{(1)} $q^\prime \cdot \Gamma \tee \lambda x.e :^{q^\prime} \,^q A \rightarrow B$.

    For the exists statement, let $\Gamma^\prime = \Gamma$.

    $q^\prime \cdot \Gamma \le q^\prime \cdot \Gamma$, as required.

    By LAM inversion on \textbf{(1)}, have \textbf{(2)} $\Gamma, x :^q A \tee e : B$.

    By LAM on \textbf{(2)}, $q_1^\prime \cdot \Gamma \tee \lambda x.e :^{q_1} \,^q A \rightarrow B$ and $q_2^\prime \cdot \Gamma \tee \lambda x.e :^{q_2} \,^q A \rightarrow B$, as required.
    \\
    \item[Case UNIT.]
    By assumption, have \textbf{(1)} $0 \cdot \Gamma \tee () :^q \unit$.

    For the exists statement, let $\Gamma^\prime = 0 \cdot \Gamma$.

    $0 \cdot \Gamma \ge q \cdot (0 \cdot \Gamma)$, as required.

    By UNIT, $q_1 \cdot (0 \cdot \Gamma) \tee () :^{q_1} \unit$ and $q_2 \cdot (0 \cdot \Gamma) \tee () :^{q_2} \unit$, as required.
    \\
    \item[Case PAIR.]
    By assumption, have \textbf{(1)} $\Gamma_1 + \Gamma_2 \tee (v_1, v_2) :^q A \times B$.

    By PAIR inversion, have \textbf{(2)} $\Gamma_1 \tee v_1 :^q A$ and \textbf{(3)} $\Gamma_2 \tee v_2 :^q B$.

    By the inductive hypothesis, have $q \cdot \Gamma_1^\prime \ge \Gamma_1$ and $q \cdot \Gamma_2^\prime \ge \Gamma_2$ where \textbf{(4)} $q_{1_1} \cdot \Gamma_1^\prime \tee v_1 :^{q_{1_1}} A$, \textbf{(5)} $q_{1_2} \cdot \Gamma_1^\prime \tee v_1 :^{q_{1_2}} A$, \textbf{(6)} $q_{2_1} \cdot \Gamma_2^\prime \tee v_2 :^{q_{2_1}} B$ and \textbf{(7)} $q_{2_2} \cdot \Gamma_2^\prime \tee v_2 :^{q_{2_2}} B$, for any $q_{1_1} + q_{1_2} = q$ and $q_{2_1} + q_{2_2} = q$.

    For the exists statement, let $\Gamma^\prime = \Gamma_1^\prime + \Gamma_2^\prime$.

    $q \cdot (\Gamma_1^\prime + \Gamma_2^\prime) \ge q \cdot \Gamma_1^\prime + q \cdot \Gamma_2^\prime \ge \Gamma_1 + \Gamma_2$, as required.

    By PAIR on \textbf{()}
    \\
    \item[Case INL.]
    By assumption, have \textbf{(1)} $\Gamma \tee \inl v :^q A + B$.

    By INL inversion, have \textbf{(2)} $\Gamma \tee v :^q A$.

    By the inductive hypothesis on \textbf{(2)}, have \textbf{(3)} $q \cdot \Gamma^\prime \le \Gamma$ where \textbf{(4)} $q_1 \cdot \Gamma^\prime \tee v :^{q_1} A$ and \textbf{(5)} $q_2 \cdot \Gamma^\prime \tee v :^{q_2} A$ for any $q_1 + q_2 = q$.

    For the exists statement, let $\Gamma^\prime = \Gamma^\prime$ as above.

    \textbf{(3)} is as required.

    By INL on each of \textbf{(4)} and \textbf{(5)}, $q_1 \cdot \Gamma^\prime \tee \inl v :^{q_1} A + B$ and $q_2 \cdot \Gamma^\prime \tee \inl v :^{q_2} A + B$, as required.
    \\
    \item[Case INR.]
    By assumption, have \textbf{(1)} $\Gamma \tee \inr v :^q A + B$.

    By INR inversion, have \textbf{(2)} $\Gamma \tee v :^q B$.

    By the inductive hypothesis on \textbf{(2)}, have \textbf{(3)} $q \cdot \Gamma^\prime \le \Gamma$ where \textbf{(4)} $q_1 \cdot \Gamma^\prime \tee v :^{q_1} B$ and \textbf{(5)} $q_2 \cdot \Gamma^\prime \tee v :^{q_2} B$ for any $q_1 + q_2 = q$.

    For the exists statement, let $\Gamma^\prime = \Gamma^\prime$ as above.

    \textbf{(3)} is as required.

    By INR on each of \textbf{(4)} and \textbf{(5)}, $q_1 \cdot \Gamma^\prime \tee \inr v :^{q_1} A + B$ and $q_2 \cdot \Gamma^\prime \tee \inr v :^{q_2} A + B$, as required.
    \\
    \item[Case $\boxempty$I.]
    By assumption, have \textbf{(1)} $q^\prime \cdot \Gamma \tee \bbox_{q^\prime} v :^q \boxempty^{q^\prime} A$.

    By $\boxempty$I inversion on \textbf{(1)}, have \textbf{(2)} $\Gamma \tee v :^{q \cdot q^\prime} A$.

    By the inductive hypothesis on \textbf{(2)}, have \textbf{(3)} $(q \cdot q^\prime) \cdot \Gamma^\prime \le \Gamma$ where \textbf{(4)} $(q_1 \cdot q^\prime) \cdot \Gamma^\prime \tee v :^{q_1 \cdot q^\prime} A$ and \textbf{(5)} $q_2 \cdot \Gamma^\prime \tee v :^{q_2 \cdot q^\prime}$ for all $q_1 + q_2 = q \cdot q^\prime$.

    For the exists statement, let $\Gamma^\prime = $
    
    By $\boxempty$I on each of \textbf{(4)} and \textbf{(5)}, $q_1 \cdot \Gamma^\prime \tee  v :^{q_1} A + B$ and $q_2 \cdot \Gamma^\prime \tee \inr v :^{q_2} A + B$, as required.
    \\
  \end{description}
\end{proof}

\begin{lemma}
\label{lem:new-modal-mode-split-2}
  If $\Gamma \tee v :^{q_1 + q_2} A$, then there exists $\Gamma_1 + \Gamma_2 \le \Gamma$ such that $\Gamma_1 \tee v :^{q_1} A$ and $\Gamma_1 \tee v :^{q_2} A$.
\end{lemma}

\begin{proof}
  Proceed by induction over the structure of the derivation of $\Gamma_1 + \Gamma_2 \tee v :^{q_1 + q_2} A$.
  \begin{description}
    \item[Case VAR.]
    By VAR, $0 \cdot \Gamma_1, x :^{q_1} A, 0 \cdot \Gamma_1 \tee x :^{q_1}$ and $0 \cdot \Gamma_2, x :^{q_2} A, 0 \cdot \Gamma_2 \tee x :^{q_2} A$ as required.
    \\
    \item[Case VSUB.]
    By assumption, have \textbf{(1)} $\Gamma, x :^{q^{\prime\prime}} A, \Gamma^\prime \tee v :^{q_1 + q_2} B$.

    By VSUB inversion on \textbf{(1)}, have \textbf{(2)} $q^\prime \le q^{\prime\prime}$ and \textbf{(3)} $\Gamma, x :^{q^\prime} A, \Gamma^\prime \tee v :^{q_1 + q_2} B$.

    By the inductive hypothesis on \textbf{(3)}, have $\Gamma_1 + \Gamma_2 \le \Gamma$, $q_1^\prime + q_2^\prime \le q^\prime$ and $\Gamma_1^\prime + \Gamma_2^\prime \le \Gamma^\prime$ such that \textbf{(4)} $\Gamma_1, x :^{q_1^\prime}, \Gamma_1^\prime \tee v :^{q_1} B$ and \textbf{(5)} $\Gamma_2, x :^{q_2^\prime}, \Gamma_2^\prime \tee v :^{q_2} B$.

    For the exists statement, let $\Gamma_1 = \Gamma_1, x :^{q_1^\prime} A, \Gamma_1^\prime$ and $\Gamma_2 = \Gamma_2, x :^{q_2^\prime} A, \Gamma_2^\prime$.

    $\Gamma_1 + \Gamma_2, x :^{q_1^\prime + q_2^\prime} A, \Gamma_1^\prime + \Gamma_2^\prime \le \Gamma, x :^{q^\prime}, \Gamma^\prime \le \Gamma, x :^{q^{\prime\prime}}, \Gamma^\prime$, as required.

    \textbf{(4)} and \textbf{(5)} are as required.
    \\
    \item[Case THUNK.]
    By assumption, have \textbf{(1)} $q \cdot \Gamma \tee \thunk e :^{q_1 + q_2} U \, B$.

    By THUNK inversion, have \textbf{(2)} $\Gamma \tee e : B$.

    For the exists statement, let $\Gamma_1 = q_1 \cdot \Gamma$ and $\Gamma_2 = q_2 \cdot \Gamma$.

    By law \ref{law:mul-distr-add-left-1}, $q_1 \cdot \Gamma + q_2 \cdot \Gamma \le (q_1 + q_2) \cdot \Gamma = q \cdot \Gamma$, as required.

    By THUNK on \textbf{(2)}, have \textbf{(3)} $q_1 \cdot \Gamma \tee \thunk e :^{q_1} U \, B$ and \textbf{(4)} $q_2 \cdot \Gamma \tee \thunk e :^{q_2} U \, B$, as required.
    \\
    \item[Case UNIT.]
    By UNIT, $0 \cdot \Gamma_1 \tee () :^{q_1} \unit$ and $0 \cdot \Gamma_1 \tee () :^{q_2} \unit$, as required.
    \\
    \item[Case PAIR.]
    By assumption, have \textbf{(1)} $\Gamma_1 + \Gamma_2 \tee (v_1, v_2) :^{q_1 + q_2} A \times B$.

    By PAIR inversion on \textbf{(1)}, have \textbf{(2)} $\Gamma_1 \tee v_1 :^{q_1 + q_2} A$ and \textbf{(3)} $\Gamma_2 \tee v_2 :^{q_1 + q_2} B$.

    By the inductive hypothesis on each of \textbf{(2)} and \textbf{(3)}, have $\Gamma_{1_1} + \Gamma_{1_2} \le \Gamma_1$ and $\Gamma_{2_1} + \Gamma_{2_2} \le \Gamma_2$ such that \textbf{(4)} $\Gamma_{1_1} \tee v_1 :^{q_1} A$, \textbf{(5)} $\Gamma_{1_2} \tee v_1 :^{q_2} A$, \textbf{(6)} $\Gamma_{2_1} \tee v_2 :^{q_1} B$ and \textbf{(7)} $\Gamma_{2_2} \tee v_2 :^{q_2} B$.

    For the exists statement, let $\Gamma_1 = \Gamma_{1_1} + \Gamma_{2_1}$ and $\Gamma_2 = \Gamma_{1_2} + \Gamma_{2_2}$.

    By \ref{law:add-mono}, $\Gamma_{1_1} + \Gamma_{2_1} + \Gamma_{1_2} + \Gamma_{2_2} \le \Gamma_1 + \Gamma_2$, as required.

    By PAIR on \textbf{(4)} with \textbf{(6)}, and \textbf{(5)} with \textbf{(7)}, $\Gamma_{1_1} + \Gamma_{2_1} \tee (v_1, v_2) :^{q_1} A \times B$ and $\Gamma_{1_2} + \Gamma_{2_2} \tee (v_1, v_2) :^{q_2} A \times B$, as required.
    \\
    \item[Case INL.]
    By assumption, have \textbf{(1)} $\Gamma \tee \inl v :^{q_1 + q_2} A + B$.

    By INL inversion on \textbf{(1)}, have \textbf{(2)} $\Gamma \tee v :^{q_1 + q_2} A$.

    By the inductive hypothesis on \textbf{(2)}, have \textbf{(3)} $\Gamma_1 + \Gamma_2 \le \Gamma$ such that \textbf{(4)} $\Gamma_1 \tee v :^{q_2} A$ and \textbf{(5)} $\Gamma_2 \tee v :^{q_2} A$.

    For the exists statement, let $\Gamma_1 = \Gamma_1$ and $\Gamma_2 = \Gamma_2$.

    \textbf{(3)} is as required.

    By INL on each of \textbf{(4)} and \textbf{(5)}, $\Gamma_1 \tee \inl v :^{q_1} A + B$ and $\Gamma_2 \tee \inl v :^{q_2} A + B$, as required.
    \\
    \item[Case INR.]
    By assumption, have \textbf{(1)} $\Gamma \tee \inr v :^{q_1 + q_2} A + B$.

    By INR inversion on \textbf{(1)}, have \textbf{(2)} $\Gamma \tee v :^{q_1 + q_2} B$.

    By the inductive hypothesis on \textbf{(2)}, have \textbf{(3)} $\Gamma_1 + \Gamma_2 \le \Gamma$ such that \textbf{(4)} $\Gamma_1 \tee v :^{q_2} B$ and \textbf{(5)} $\Gamma_2 \tee v :^{q_2} B$.

    For the exists statement, let $\Gamma_1 = \Gamma_1$ and $\Gamma_2 = \Gamma_2$.

    \textbf{(3)} is as required.

    By INR on each of \textbf{(4)} and \textbf{(5)}, $\Gamma_1 \tee \inr v :^{q_1} A + B$ and $\Gamma_2 \tee \inl v :^{q_2} A + B$, as required.
    \\
    \item[Case $\boxempty$I.] 
    By assumption, have \textbf{(1)} $\Gamma \tee \bbox_q v :^{q_1 + q_2} \boxempty^q A$.

    By $\boxempty$I inversion on \textbf{(1)}, have \textbf{(2)} $\Gamma \tee v :^{(q_1 + q_2) \cdot q} A$.

    By weakening using law \ref{law:mul-distr-add-right-2} on \textbf{(2)}, have \textbf{(3)} $\Gamma \tee v :^{q_1 \cdot q + q_2 \cdot q} A$.

    By the inductive hypothesis on \textbf{(3)}, have \textbf{(4)} $\Gamma_1 + \Gamma_2 \le \Gamma$ such that \textbf{(5)} $\Gamma_1 \tee v :^{q_1 \cdot q} A$ and \textbf{(6)} $\Gamma_2 \tee v :^{q_2 \cdot q} A$.

    For the exists statement, let $\Gamma_1 = \Gamma_1$ and $\Gamma_2 = \Gamma_2$.

    \textbf{(4)} is as required.

    By $\boxempty$I on each of \textbf{(4)} and \textbf{(5)}, $\Gamma_1 \tee \bbox_q v :^{q_1} \boxempty^q A$ and $\Gamma_2 \tee \bbox_q v :^{q_2} \boxempty^q A$, as required.
  \end{description}
\end{proof}

\begin{lemma}
\label{lem:new-modal-div-both-sides}
  If $\Gamma \tee v :^{q^\prime \cdot q} A$, then $\Gamma \setminus q^\prime \tee v :^q A$.
\end{lemma}

\begin{proof}
  Proceed by induction over the structure of the derivation of $\Gamma \tee v :^{q^\prime \cdot q} A$.
  \begin{description}
    \item[Case VAR ???.]
    By VAR, have \textbf{(1)} $(0 \cdot \Gamma) \setminus q^\prime, x :^q A, (0 \cdot \Gamma^\prime) \setminus q^\prime \tee x :^q A$. ???

    By weakening using \ref{lem:mul-left-inverse-div-le}, $(0 \cdot \Gamma) \setminus q^\prime, x :^{(q^\prime \cdot q) \setminus q^\prime} A, (0 \cdot \Gamma^\prime) \setminus q^\prime \tee x :^q A$, as required.
    \\
    \item[Case VSUB.]
    By assumption, have \textbf{(1)} $\Gamma, x :^{q_2} A, \Gamma^\prime \tee v :^{q^\prime \cdot q} B$. 

    By VSUB inversion on \textbf{(1)}, have \textbf{(2)} $q_1 \le q_2$ and \textbf{(3)} $\Gamma, x :^{q_1}, \Gamma^\prime \tee v :^{q^\prime \cdot q} B$.

    By the inductive hypothesis on \textbf{(3)}, have \textbf{(4)} $\Gamma \setminus q^\prime, x :^{q_1 \setminus q^\prime} A, \Gamma^\prime \setminus q^\prime \tee v :^q B$.

    By \ref{lem:div-mono} on \textbf{(2)}, have \textbf{(5)} $q_1 \setminus q^\prime \le q_2 \setminus q^\prime$.

    By VSUB on \textbf{(4)} and \textbf{(5)}, $\Gamma \setminus q^\prime, x :^{q_2 \setminus q^\prime} A, \Gamma^\prime \setminus q^\prime \tee v :^q B$, as required.
    \\
    \item[Case THUNK.] 
    By assumption, have \textbf{(1)} $(q^\prime \cdot q) \cdot \Gamma \tee \thunk e :^{q^\prime \cdot q} U \, \underline{B}$.

    By THUNK inversion on \textbf{(1)}, have \textbf{(2)} $\Gamma \tee e : \underline{B}$.

    By THUNK on \textbf{(2)}, have \textbf{(3)} $q \cdot \Gamma \tee \thunk e :^q U \, \underline{B}$.

    By weakening using law \ref{lem:mul-left-inverse-div-le} on \textbf{(3)}, have \textbf{(4)} $(q^\prime \cdot (q \cdot \Gamma)) \setminus q^\prime \tee \thunk e :^q U \, \underline{B}$.

    By weakening using \ref{law:mul-assoc-1} on \textbf{(4)}, have $((q^\prime \cdot q) \cdot \Gamma) \setminus q^\prime \tee \thunk e :^q U \, \underline{B}$, as required.
    \\
    \item[Case UNIT ???.]
    By UNIT, have $0 \cdot (\Gamma \setminus q^\prime) \tee () :^q \unit$.
    \\
    \item[Case PAIR.]
    By assumption, have \textbf{(1)} $\Gamma_1 + \Gamma_2 \tee (v_1, v_2) :^{q^\prime \cdot q} A \times B$.

    By PAIR inversion on \textbf{(1)}, have \textbf{(2)} $\Gamma_1 \tee v_1 :^{q^\prime \cdot q} A$ and \textbf{(3)} $\Gamma_2 \tee v_2 :^{q^\prime \cdot q} B$.

    By the inductive hypothesis on each of \textbf{(2)} and \textbf{(3)}, have \textbf{(4)} $\Gamma_1 \setminus q^\prime \tee v_1 :^q A$ and \textbf{(5)} $\Gamma_2 \setminus q^\prime \tee v_2 :^q B$.

    By PAIR on \textbf{(4)} and \textbf{(5)}, have \textbf{(6)} $\Gamma_1 \setminus q^\prime + \Gamma_2 \setminus q^\prime \tee (v_1, v_2) :^q A \times B$.

    By \ref{lem:div-distr-add} on \textbf{(6)}, $(\Gamma_1 + \Gamma_2) \setminus q^\prime \tee (v_1, v_2) :^q A \times B$, as required.
    \\
    \item[Case INL.]
    By assumption, have \textbf{(1)} $\Gamma \tee \inl v :^{q^\prime \cdot q} A + B$.

    By INL inversion on \textbf{(1)}, have \textbf{(2)} $\Gamma \tee v :^{q^\prime \cdot q} A$.

    By the inductive hypothesis on \textbf{(2)}, have \textbf{(3)} $\Gamma \setminus q^\prime \tee v :^q A$

    By INL on \textbf{(3)}, $\Gamma \setminus q^\prime \tee : \inl v :^q A + B$, as required.
    \\
    \item[Case INR.]
    By assumption, have \textbf{(1)} $\Gamma \tee \inr v :^{q^\prime \cdot q} A + B$.

    By INR inversion on \textbf{(1)}, have \textbf{(2)} $\Gamma \tee v :^{q^\prime \cdot q} B$.

    By the inductive hypothesis on \textbf{(2)}, have \textbf{(3)} $\Gamma \setminus q^\prime \tee v :^q B$

    By INR on \textbf{(3)}, $\Gamma \setminus q^\prime \tee : \inr v :^q A + B$, as required.
    \\
    \item[Case $\boxempty$I TODO.]
    By assumption, have \textbf{(1)} $\Gamma \tee \bbox_q v :^{q_1 \cdot q_2} \boxempty^q A$.

    By $\boxempty$I inversion on \textbf{(1)}, have \textbf{(2)} $\Gamma \tee v :^{(q_1 \cdot q_2) \cdot q} A$.

    By weakening using \ref{law:mul-assoc-2} on \textbf{(2)}, have \textbf{(3)} $\Gamma \tee v :^{q_1 \cdot (q_2 \cdot q)} A$. ???

    By the inductive hypothesis on \textbf{(3)}, have \textbf{(4)} $\Gamma \setminus q_1 \tee v :^{q_2 \cdot q} A$.

    By $\boxempty$I on \textbf{(4)}, $\Gamma \setminus q_1 \tee \bbox_q v :^{q_2} \boxempty^q A$, as required.
  \end{description}
\end{proof}

\begin{lemma}
\label{lem:new-modal-invariant}
  If $\Gamma \tee v :^q A$, then there exists a $\Gamma^\prime$ such that $q \cdot \Gamma^\prime \le \Gamma$ and $q \cdot \Gamma^\prime \tee v :^q A$.
\end{lemma}

\begin{proof}
  TODO.
\end{proof}

\begin{lemma}[Forward Substitution for New Type Theory Values]
\label{lem:new-modal-forward-substitution-vals}
  If $\Gamma_1 \tee v :^q A$ and $\Gamma_2, x :^q A, \Gamma_2^\prime \tee v^\prime :^{q^\prime} B$, then $\Gamma_1 + \Gamma_2, \Gamma_2^\prime \tee v^\prime[v/x] :^{q^\prime} B$.
\end{lemma}

\begin{proof}
  Proceed by induction over the structure of the derivation of $\Gamma_2, x :^q A, \Gamma_2^\prime \tee v^\prime :^{q^\prime} B$.

  \begin{description}
    \item[Case VAR ???.]
    Since $\Gamma_2 = \Gamma_2^\prime = \cdot$, by assumption, $\Gamma_1 \tee v : A$, as required.
    \\
    \item[Case VSUB.]
    By assumption, have \textbf{(1)} $\Gamma_1 \tee v :^q A$ and \textbf{(2)} $\Gamma_2, x :^q A, y :^{q_2} B, \Gamma_2^\prime \tee v^\prime :^{q^\prime} C$.

    By VSUB inversion on \textbf{(2)}, have \textbf{(3)} $q_1 \le q_2$ and \textbf{(4)} $\Gamma_2, x :^q A, y :^{q_1} B, \Gamma_2^\prime \tee v^\prime :^{q^\prime} C$.

    By the inductive hypothesis on \textbf{(1)} and \textbf{(3)}, have \textbf{(5)} $\Gamma_1 + \Gamma_2, y :^{q_1} B, \Gamma_2^\prime \tee v^\prime[v/x] :^{q^\prime} C$.

    By VSUB on \textbf{(3)} and \textbf{(5)}, $\Gamma_1 + \Gamma_2, y :^{q_2} B, \Gamma_2^\prime \tee v^\prime[v/x] :^{q^\prime} C$, as required.
    \\
    \item[Case THUNK.]
    By assumption, have \textbf{(1)} $\Gamma_1 \tee v :^{q^\prime \cdot q} A$ and \textbf{(2)} $q^\prime \cdot \Gamma_2, x :^{q^\prime \cdot q} A, \Gamma_2^\prime \tee \thunk e :^{q^\prime} U \, \underline{B}$.

    By THUNK inversion on \textbf{(2)}, have \textbf{(3)} $\Gamma_2, x :^q A, \Gamma_2^\prime \tee e : \underline{B}$.

    By \ref{lem:new-modal-div-both-sides} on \textbf{(1)}, have \textbf{(4)} $\Gamma_1 \setminus q^\prime \tee v :^q A$.

    By \ref{lem:new-modal-forward-substitution-exprs} on \textbf{(3)} and \textbf{(4)}, have \textbf{(5)} $\Gamma_1 \setminus q^\prime + \Gamma_2, \Gamma_2^\prime \tee e[v/x] : \underline{B}$.

    By THUNK on \textbf{(3)}, have \textbf{(6)} $q^\prime \cdot (\Gamma_1 \setminus q^\prime + \Gamma_2), q^\prime \cdot \Gamma_2^\prime \tee \thunk e[v/x] :^{q^\prime} \underline{B}$.

    By weakening using \ref{law:mul-distr-add-left-2} on \textbf{(6)}, have \textbf{(7)} $q^\prime \cdot (\Gamma_1 \setminus q^\prime) + q^\prime \cdot \Gamma_2, q^\prime \cdot \Gamma_2^\prime \tee \thunk e[v/x] :^{q^\prime} \underline{B}$.

    By weakening using \ref{law:mul-inverse-div} on \textbf{(7)}, $\Gamma_1 + q^\prime \cdot \Gamma_2, q^\prime \cdot \Gamma_2^\prime \tee \thunk e[v/x] :^{q^\prime} \underline{B}$, as required.
    \\
    \item[Case UNIT.]
    By assumption, have \textbf{(1)} $\Gamma_1 \tee v :^0 A$ and \textbf{(2)} $0 \cdot \Gamma_2, x :^0 A, 0 \cdot \Gamma_2^\prime \tee () :^q \unit$.

    By \ref{lem:new-modal-invariant} on \textbf{(1)}, have some $\Gamma_1^\prime$ such that \textbf{(3)} $0 \cdot \Gamma_1^\prime \le \Gamma_1$ and \textbf{(4)} $0 \cdot \Gamma_1^\prime \tee v :^0 A$.

    By UNIT, have \textbf{(5)} $0 \cdot (\Gamma_1^\prime + \Gamma_2), 0 \cdot \Gamma_2^\prime \tee () :^q \unit$.

    By weakening using law \ref{law:mul-distr-add-left-2} on \textbf{(5)}, have \textbf{(6)} $0 \cdot \Gamma_1^\prime + 0 \cdot \Gamma_2, 0 \cdot \Gamma_2^\prime \tee () :^q \unit$.

    By weakening using \textbf{(3)} on \textbf{(6)}, $\Gamma_1 + 0 \cdot \Gamma_2, 0 \cdot \Gamma_2^\prime \tee () :^q \unit$, as required.
    \\
    \item[Case PAIR.]
    By assumption, have \textbf{(1)} $\Gamma_1 \tee v :^{q_1 + q_2} A$ and \textbf{(2)} $\Gamma_{2_1} + \Gamma_{2_2}, x :^{q_1 + q_2} A, \Gamma_{2_1}^\prime + \Gamma_{2_2}^\prime \tee (v_1, v_2) :^{q^\prime} B_1 \times B_2$. 

    By PAIR inversion on \textbf{(2)}, have \textbf{(3)} $\Gamma_{2_1}, x :^{q_1} A, \Gamma_{2_1}^\prime \tee v_1 :^{q^\prime} B_1$ and \textbf{(4)} $\Gamma_{2_2}, x :^{q_1} A, \Gamma_{2_2}^\prime \tee v_2 :^{q^\prime} B_2$.

    By \ref{lem:new-modal-mode-split-2} on \textbf{(1)}, have \textbf{(5)} $\Gamma_{1_1} + \Gamma_{1_2} \le \Gamma_1$ such that \textbf{(6)} $\Gamma_{1_1} \tee v :^{q_1} A$ and \textbf{(7)} $\Gamma_{1_2} \tee v :^{q_2} A$.

    By the inductive hypothesis on \textbf{(3)} with \textbf{(6)} and \textbf{(4)} with \textbf{(7)}, have \textbf{(8)} $\Gamma_{1_1} + \Gamma_{2_1}, \Gamma_{2_1}^\prime \tee v_1[v/x] :^{q^\prime} B_1$ and \textbf{(9)} $\Gamma_{1_2} + \Gamma_{2_2}, \Gamma_{2_2}^\prime \tee v_2[v/x] :^{q^\prime} B_2$.

    By PAIR on \textbf{(8)} and \textbf{(9)}, $\Gamma_1 + \Gamma_2, \Gamma_2^\prime \tee (v_1[v/x], v_2[v/x]) :^{q^\prime} B_1 \times B_2$, as required.
    \\
    \item[Case INL.]
    By assumption, have \textbf{(1)} $\Gamma_1 \tee v :^q A$ and \textbf{(2)} $\Gamma_2, x :^q A, \Gamma_2^\prime \tee \inl v^\prime :^{q^\prime} B_1 + B_2$.

    By INL inversion on \textbf{(2)}, have \textbf{(3)} $\Gamma_2, x :^q A, \Gamma_2^\prime \tee v^\prime :^{q^\prime} B_1$.

    By the inductive hypothesis on \textbf{(1)} and \textbf{(3)}, have \textbf{(4)} $\Gamma_1 + \Gamma_2, \Gamma_2^\prime \tee v^\prime[v/x] :^{q^\prime} B_1$.

    By INL on \textbf{(4)}, $\Gamma_1 + \Gamma_2, \Gamma_2^\prime \tee \inl v^\prime[v/x] : B_1 + B_2$, as required.
    \\
    \item[Case INR.]
    By assumption, have \textbf{(1)} $\Gamma_1 \tee v :^q A$ and \textbf{(2)} $\Gamma_2, x :^q A, \Gamma_2^\prime \tee \inr v^\prime :^{q^\prime} B_1 + B_2$.

    By INR inversion on \textbf{(2)}, have \textbf{(3)} $\Gamma_2, x :^q A, \Gamma_2^\prime \tee v^\prime :^{q^\prime} B_2$.

    By the inductive hypothesis on \textbf{(1)} and \textbf{(3)}, have \textbf{(4)} $\Gamma_1 + \Gamma_2, \Gamma_2^\prime \tee v^\prime[v/x] :^{q^\prime} B_2$.

    By INR on \textbf{(4)}, $\Gamma_1 + \Gamma_2, \Gamma_2^\prime \tee \inr v^\prime[v/x] : B_1 + B_2$, as required.
    \\
    \item[Case $\boxempty$I.]
    By assumption, have \textbf{(1)} $\Gamma_1 \tee v :^q A$ and \textbf{(2)} $\Gamma_2, x :^q A, \Gamma_2^\prime \tee \bbox_{q_2} v^\prime :^{q_1} \boxempty^{q_2} A$.

    By $\boxempty$I inversion on \textbf{(2)}, have \textbf{(3)} $\Gamma_2, x :^q A, \Gamma_2^\prime \tee v^\prime :^{q_1 \cdot q_2} A$.

    By the inductive hypothesis on \textbf{(1)} and \textbf{(3)}, have \textbf{(4)} $\Gamma_1 + \Gamma_2, \Gamma_2^\prime \tee v^\prime[v/x] :^{q_1 \cdot q_2} A$.

    By $\boxempty$I on \textbf{(4)}, $\Gamma_1 + \Gamma_2, \Gamma_2^\prime \tee \bbox_{q_2} v^\prime :^{q_1} \boxempty_{q_2} A$, as required.
  \end{description}
\end{proof}

\begin{lemma}[Forward Substitution for New Type Theory Expressions]
\label{lem:new-modal-forward-substitution-exprs}
  If $\Gamma_1 \tee v :^q A$ and $\Gamma_2, x :^q A, \Gamma_2^\prime \tee e : B$, then $\Gamma_1 + \Gamma_2, \Gamma_2^\prime \tee e[v/x] : B$.
\end{lemma}

\begin{proof}
  Proceed by induction over the structure of the derivation of $\Gamma_2, x :^q A, \Gamma_2^\prime \tee e : B$.
  
  \begin{description}
    \item[Case RETURN.]
    By assumption, have \textbf{(1)} $\Gamma_1 \tee v :^q A$ and $\Gamma_2, x :^q A, \Gamma_2^\prime \tee \return v^\prime : F^{q^\prime} A$.

    By RETURN inversion on \textbf{(2)}, have \textbf{(3)} $\Gamma_2, x :^q A, \Gamma_2^\prime \tee v :^{q^\prime} A$.

    By \ref{lem:new-modal-forward-substitution-vals} on \textbf{(1)} and \textbf{(3)}, have \textbf{(4)} $\Gamma_1 + \Gamma_2, \Gamma_2^\prime \tee v^\prime[v/x] :^{q^\prime} A$.

    Thus, by RETURN on \textbf{(4)}, $\Gamma_1 + \Gamma_2, \Gamma_2^\prime \tee \return v^\prime[v/x] : F^{q^\prime} A$, as required.
    \\
    \item[Case FORCE.]
    By assumption, have \textbf{(1)} $\Gamma_1 \tee v :^q A$ and \textbf{(2)} $\Gamma_2, x :^q A, \Gamma_2^\prime \tee \force v^\prime : \underline{B}$.

    By FORCE inversion on \textbf{(2)}, have \textbf{(3)} $\Gamma_2, x :^q A, \Gamma_2^\prime \tee v^\prime :^1 U \, \underline{B}$.

    By \ref{lem:new-modal-forward-substitution-vals} on \textbf{(1)} and \textbf{(3)}, have \textbf{(4)} $\Gamma_1 + \Gamma_2, \Gamma_2^\prime \tee v^\prime[v/x] :^1 U \, \underline{B}$.

    By FORCE on \textbf{(4)}, $\Gamma_1 + \Gamma_2, \Gamma_2^\prime \tee \force v^\prime[v/x] : \underline{B}$, as required.
    \\
    \item[Case ESUB.]
    By assumption, have \textbf{(1)} $\Gamma_1 \tee v :^q A$ and \textbf{(2)} $\Gamma_2, x :^q A, y :^{q_2} B, \Gamma_2^\prime \tee e : C$.

    By ESUB inversion on \textbf{(2)}, have \textbf{(3)} $q_1 \le q_2$ and \textbf{(4)} $\Gamma_2, x :^q A, y :^{q_1} B, \Gamma_2^\prime \tee e : C$.

    By the inductive hypothesis on \textbf{(1)} and \textbf{(3)}, have \textbf{(5)} $\Gamma_1 + \Gamma_2, y :^{q_1} B, \Gamma_2^\prime \tee e[v/x] : C$.

    By ESUB on \textbf{(3)} and \textbf{(5)}, $\Gamma_1 + \Gamma_2, y :^{q_2} B, \Gamma_2^\prime \tee e[v/x] : C$, as required.
    \\
    \item[Case LAM.]
    By assumption, have \textbf{(1)} $\Gamma_1 \tee v :^q A$ and \textbf{(2)} $\Gamma_2, x :^q A, \Gamma_2^\prime \tee \lambda y.e : \,^{q^\prime}B_1 \rightarrow \underline{B_2}$.

    By LAM inversion on \textbf{(2)}, have \textbf{(3)} $\Gamma_2, x :^q A, \Gamma_2^\prime, y :^{q^\prime} B_1 \tee e : \underline{B_2}$.

    By the inductive hypothesis on \textbf{(3)}, have \textbf{(4)} $\Gamma_1 + \Gamma_2, \Gamma_2^\prime, y :^{q^\prime} B_1 \tee e[v/x] : \underline{B_2}$.

    By LAM on \textbf{(4)}, $\Gamma_1 + \Gamma_2, \Gamma_2^\prime \tee \lambda y.e :^{q^\prime} B_1 \rightarrow \underline{B_2}$, as required.
    \\
    \item[Case APP.]
    By assumption, have \textbf{(1)} $\Gamma_1 \tee v :^{q_1 + q_2} A$ and \textbf{(2)} $\Gamma_{2_1} + \Gamma_{2_2}, x :^{q_1 + q_2} A, \Gamma_{2_1}^\prime + \Gamma_{2_2}^\prime \tee e \, v : \underline{B}$.

    By APP inversion on \textbf{(2)}, have \textbf{(3)} $\Gamma_{2_1}, x :^{q_1}, \Gamma_{2_1}^\prime \tee e : \,^{q^\prime} B_1 \rightarrow \underline{B_2}$ and \textbf{(4)} $\Gamma_{2_2}, x :^{q_2}, \Gamma_{2_2}^\prime \tee v :^{q^\prime} B_1$.

    By \ref{lem:new-modal-mode-split-2} on \textbf{(1)}, have some \textbf{(5)} $\Gamma_{1_1} + \Gamma_{1_2} \le \Gamma_1$ such that \textbf{(6)} $\Gamma_{1_1} \tee v :^{q_1^\prime} A$ and \textbf{(7)} $\Gamma_{1_2} \tee v :^{q_2^\prime} A$.
    
    By the inductive hypothesis on \textbf{(3)} and \textbf{(6)}, have \textbf{(8)} $\Gamma_{1_1} + \Gamma_{2_1}, \Gamma_{2_1}^\prime \tee e[v/x] : \,^{q^\prime} B_1 \rightarrow \underline{B_2}$. 

    By \ref{lem:new-modal-forward-substitution-vals} on \textbf{(4)} and \textbf{(7)}, have \textbf{(9)} $\Gamma_{1_2} + \Gamma_{2_2}, \Gamma_{2_2}^\prime \tee v^\prime[v/x] :^{q^\prime} B_1$.

    By APP on \textbf{(8)} and \textbf{(9)}, have \textbf{(10)} $\Gamma_{1_1} + \Gamma_{1_2} + \Gamma_2, \Gamma_2^\prime \tee e[v/x] \, v^\prime[v/x] : \underline{B_2}$.

    By weakening on \textbf{(10)} using \textbf{(5)}, $\Gamma_1 + \Gamma_2, \Gamma_2^\prime \tee e[v/x] \, v^\prime[v/x] : \underline{B_2}$, as required.
    \\
    \item[Case BIND.]
    By assumption, have \textbf{(1)} $\Gamma_1 \tee v :^{q_1 + q_2} A$ and \textbf{(2)} $\Gamma_{2_1} + \Gamma_{2_2}, x :^{q_1 + q_2}, \Gamma_{2_1}^\prime + \Gamma_{2_2}^\prime \tee \llet y \leftarrow e_1 \lin e_2 : \underline{B_2}$.

    By BIND inversion on \textbf{(2)}, have \textbf{(3)} $\Gamma_{2_1}, x :^{q_1} A, \Gamma_{2_1}^\prime \tee e_1 : F^{q^\prime} B_1$ and \textbf{(4)} $\Gamma_{2_2}, x :^{q_2} A, \Gamma_{2_2}^\prime, y :^{q^\prime} B_1 \tee e_2 : \underline{B_2}$.
    
    By \ref{lem:new-modal-mode-split-2} on \textbf{(1)}, have some \textbf{(5)} $\Gamma_{1_1} + \Gamma_{1_2} \le \Gamma_1$ such that \textbf{(6)} $\Gamma_{1_1} \tee v :^{q_1^\prime} A$ and \textbf{(7)} $\Gamma_{1_2} \tee v :^{q_2^\prime} A$.

    By the inductive hypothesis on \textbf{(3)} with \textbf{(6)} and \textbf{(4)} with \textbf{(7)}, have \textbf{(8)} $\Gamma_{1_1} + \Gamma_{2_1}, \Gamma_{2_1}^\prime \tee e_1[v/x] : F^{q^\prime} B_1$ and \textbf{(9)} $\Gamma_{1_2} + \Gamma_{2_2}, \Gamma_{2_2}^\prime, y :^{q^\prime} B_1 \tee e_2[v/x] : \underline{B_2}$.

    By BIND on \textbf{(8)} and \textbf{(9)}, have \textbf{(10)} $\Gamma_{1_1} + \Gamma_{1_2} + \Gamma_2, \Gamma_2^\prime \tee \llet y \leftarrow e_1 \lin e_2 : \underline{B_2}$.

    By weakening on \textbf{(10)} using \textbf{(5)}, $\Gamma_1 + \Gamma_2, \Gamma_2^\prime \tee \llet y \leftarrow e_1 \lin e_2 : \underline{B_2}$, as required.
    \\
    \item[Case $\boxempty$E.]
    By assumption, have \textbf{(1)} $\Gamma_1 \tee v :^{q_1 + q_2} A$ and \textbf{(2)} $\Gamma_{2_1} + \Gamma_{2_2}, x :^{q_1^\prime + q_2^\prime} A, \Gamma_{2_1}^\prime + \Gamma_{2_2}^\prime \tee \llet_{q_1} \bbox_{q_2} y = \bbox_{q_2} v^\prime \lin e : \underline{B_2}$.

    By $\boxempty$E inversion on \textbf{(2)}, have \textbf{(3)} $\Gamma_{2_1}, x :^{q_1^\prime} A, \Gamma_{2_1}^\prime \tee \bbox_{q_2} v^\prime :^{q_1} \boxempty^{q_2} B_1$ and \textbf{(4)} $\Gamma_{2_2}, x :^{q_2^\prime} A, \Gamma_{2_2}^\prime, y :^{q_1 \cdot q_2} B_1 \tee e : \underline{B_2}$.

    By \ref{lem:new-modal-mode-split-2} on \textbf{(1)}, have some \textbf{(5)} $\Gamma_{1_1} + \Gamma_{1_2} \le \Gamma_1$ such that \textbf{(6)} $\Gamma_{1_1} \tee v :^{q_1^\prime} A$ and \textbf{(7)} $\Gamma_{1_2} \tee v :^{q_2^\prime} A$.

    By \ref{lem:new-modal-forward-substitution-vals} on \textbf{(3)} and \textbf{(6)}, have \textbf{(8)} $\Gamma_{1_1} + \Gamma_{2_1}, \Gamma_{2_1}^\prime \tee \bbox_{q_2} v^\prime[v/x] :^{q_1} \boxempty^{q_2} B_1$. 

    By the inductive hypothesis on \textbf{(4)} and \textbf{(7)}, have \textbf{(9)} $\Gamma_{1_2} + \Gamma_{2_2}, \Gamma_{2_2}^\prime, y :^{q_1 \cdot q_2} B_1 \tee e[v/x] : B_2$.

    By $\boxempty$E on \textbf{(8)} and \textbf{(9)}, have \textbf{(10)} $\Gamma_{1_1} + \Gamma_{1_2} + \Gamma_2, \Gamma_2^\prime \tee \llet_{q_1} \bbox_{q_2} y = \bbox_{q_2} v^\prime[v/x] \lin e[v/x] : B_2$.

    By weakening on \textbf{(10)} using \textbf{(5)}, $\Gamma_1 + \Gamma_2, \Gamma_2^\prime \tee \llet_{q_1} \bbox_{q_2} y = \bbox_{q_2} v^\prime[v/x] \lin e[v/x] : B_2$, as required.
    \\
    \item[Case SPLIT.]
    By assumption, have \textbf{(1)} $\Gamma_1 \tee v :^{q_1 + q_2} A$ and $\Gamma_{2_1} + \Gamma_{2_2}, x^{q_1 + q_2} A, \Gamma_{2_2} + \Gamma_{2_2}^\prime \tee \llet_{q^\prime} (y, z) = v^\prime \lin e : \underline{C}$.

    By SPLIT inversion on \textbf{(2)}, have \textbf{(3)} $\Gamma_{2_1}, x :^{q_1}, \Gamma_{2_2} \tee v :^{q^\prime} B_1 \times B_2$ and \textbf{(4)} $\Gamma_{2_2}, x :^{q_2}, \Gamma_{2_2}^\prime, y :^{q^\prime} B_1, z :^{q^\prime} B_2 \tee e : \underline{C}$.

    By \ref{lem:new-modal-mode-split-2} on \textbf{(1)}, have some \textbf{(5)} $\Gamma_{1_1} + \Gamma_{1_2} \le \Gamma_1$ such that \textbf{(6)} $\Gamma_{1_1} \tee v :^{q_1^\prime} A$ and \textbf{(7)} $\Gamma_{1_2} \tee v :^{q_2^\prime} A$.
    
    By \ref{lem:new-modal-forward-substitution-vals} on \textbf{(3)} and \textbf{(6)}, have \textbf{(8)} $\Gamma_{1_1} + \Gamma_{2_1}, \Gamma_{2_1}^\prime \tee v^\prime[v/x] :^{q^\prime} B_1 \times B_2$. 

    By the inductive hypothesis on \textbf{(4)} and \textbf{(7)}, have \textbf{(9)} $\Gamma_{1_2} + \Gamma_{2_2}, \Gamma_{2_2}^\prime, y :^{q^\prime} B_1, z :^{q^\prime} B_2 \tee e[v/x] : \underline{C}$.

    By SPLIT on \textbf{(8)} and \textbf{(9)}, have \textbf{(10)} $\Gamma_{1_1} + \Gamma_{1_2} + \Gamma_2, \Gamma_2^\prime \tee \llet_{q^\prime} (y, z) = v^\prime[v/x] \lin e[v/x] : \underline{C}$.

    By weakening on \textbf{(10)} using \textbf{(5)}, $\Gamma_1 + \Gamma_2, \Gamma_2^\prime \tee \llet_{q^\prime} (y, z) = v^\prime[v/x] \lin e[v/x] : \underline{C}$, as required.
    \\
    \item[Case CASE.]
    By assumption, have \textbf{(1)} $\Gamma_1 \tee v :^{q_1 + q_2} A$ and $\Gamma_{2_1} + \Gamma_{2_2}, x :^{q_1 + q_2} A, \Gamma_{2_1}^\prime + \Gamma_{2_2}^\prime \tee \case_{q^\prime} \, v^\prime \{y.e_1; \, z.e_2\} : \underline{C}$.

    By CASE inversion on \textbf{(2)}, have \textbf{(3)} $\Gamma_{2_1}, x :^{q_1} A, \Gamma_{2_1}^\prime \tee v^\prime :^{q^\prime} B_1 + B_2$, \textbf{(4)} $\Gamma_{2_2}, x :^{q_2} A, \Gamma_{2_2}^\prime, y :^{q^\prime} B_1 \tee e_1 : \underline{C}$ and \textbf{(5)} $\Gamma_{2_2}, x :^{q_2} A, \Gamma_{2_2}^\prime, z :^{q^\prime} B_2 \tee e_2 : \underline{C}$.

    By \ref{lem:new-modal-mode-split-2} on \textbf{(1)}, have some \textbf{(6)} $\Gamma_{1_1} + \Gamma_{1_2} \le \Gamma_1$ such that \textbf{(7)} $\Gamma_{1_1} \tee v :^{q_1^\prime} A$ and \textbf{(8)} $\Gamma_{1_2} \tee v :^{q_2^\prime} A$.
    
    By \ref{lem:new-modal-forward-substitution-vals} on \textbf{(3)} and \textbf{(7)}, have \textbf{(9)} $\Gamma_{1_1} + \Gamma_{2_1}, \Gamma_{2_1}^\prime \tee v^\prime[v/x] :^{q^\prime} B_1 + B_2$. 

    By the inductive hypothesis on each of \textbf{(4)} and \textbf{(5)} with \textbf{(8)}, have \textbf{(10)} $\Gamma_{1_2} + \Gamma_{2_2}, \Gamma_{2_2}^\prime, y :^{q^\prime} B_1 \tee e_1[v/x] : \underline{C}$ and \textbf{(11)} $\Gamma_{1_2} + \Gamma_{2_2}, \Gamma_{2_2}^\prime, z :^{q^\prime} B_2 \tee e_2[v/x] : \underline{C}$.

    By CASE on \textbf{(9)}, \textbf{(10)}, \textbf{(11)}, have \textbf{(12)} $\Gamma_{1_1} + \Gamma_{1_2} + \Gamma_2, \Gamma_2^\prime \tee \case_{q^\prime} \, v^\prime[v/x] \{y.e_1[v/x]; \, z.e_2[v/x]\} : \underline{C}$.
    
    By weakening on \textbf{(12)} using \textbf{(6)}, $\Gamma_1 + \Gamma_2, \Gamma_2^\prime \tee \case_{q^\prime} \, v^\prime[v/x] \{y.e_1[v/x]; \, z.e_2[v/x]\}$, as required.
    \\
    \item[Case ABSURD.]
    By assumption, have \textbf{(1)} $\Gamma_1 \tee v :^\top A$ and $\top \cdot \Gamma_2, x :^\top A, \top \cdot \Gamma_2^\prime \tee \absurd e : \underline{B}$.

    By ABSURD inversion on \textbf{(2)}, have \textbf{(3)} $\Gamma_2, x :^q A, \Gamma_2^\prime \tee e : \zero$.

    By \ref{lem:new-modal-div-both-sides} on \textbf{(1)}, have \textbf{(4)} $\bot \cdot \Gamma_1 \tee v :^q A$.

    By the inductive hypothesis on \textbf{(3)} and \textbf{(4)}, have \textbf{(5)} $\bot \cdot \Gamma_1 + \Gamma_2, \Gamma_2^\prime \tee e : \zero$.

    By ABSURD on \textbf{(5)}, have \textbf{(6)} $\top \cdot (\bot \cdot \Gamma_1 + \Gamma_2), \top \cdot \Gamma_2^\prime \tee \absurd e : \underline{B}$.

    By weakening on \textbf{(6)} using law \ref{law:mul-distr-add-left-2}, have \textbf{(7)} $\top \cdot \Gamma_1 + \top \cdot \Gamma_2, \top \cdot \Gamma_2^\prime \tee \absurd e : \underline{B}$.
    
    By weakening on \textbf{(6)} considering that $\top$ is the least element, $\Gamma_1 + \top \cdot \Gamma_2, \top \cdot \Gamma_2^\prime \tee \absurd e : \underline{B}$, as required.
  \end{description}
\end{proof}

Assume that $\Gamma \tee \llet_{q_1} \bbox_{q_2} x = \bbox_{q_2} v \lin e : \underline{B}$.
By the box-elim rule, we have $\Gamma = \Gamma_1 + \Gamma_2$ for some $\Gamma_1$ and $\Gamma_2$,
and $\Gamma_1 \tee \bbox_{q_2} v :^{q_1} \boxempty^{q_2} A$,
and $\Gamma_2, x :^{q_1 \cdot q_2} A \tee e : \underline{B}$.
By the box-intro rule, we have $\Gamma_1 \tee v :^{q_1 \cdot q_2} A$.
We use the forward substitution lemma, to obtain $\Gamma_1 + \Gamma_2 \tee e[v/x] : \underline{B}$.

\begin{acks}

The internship was supervised by Anton Lorenzen and Sam Lindley, University of Edinburgh. \\
The QED square was provided by Bharat Haria, University of Warwick.

\end{acks}

\bibliographystyle{ACM-Reference-Format}
\bibliography{references.bib}

\end{document}